% Sección: Introducción

El objetivo del trabajo práctico fue construir dos clasificadores que, a partir
de cierta información acerca de las características de las células que componen
un tumor, permitan identificar si este es maligno o benigno con un grado
aceptable de error bajo algún criterio establecido. Para ello se contó con una
muestra de observaciones de tumores de los que se registraron tres medidas
resumen de diez características de una cantidad desconocida de sus células
mediante imágenes microscópicas \parencite{williamwolberg_1993_breast}. Los
clasificadores obtenidos se seleccionaron de un total de seis que fueron
entrenados y evaluados sobre los datos disponibles según su desempeño en
ciertas métricas que reflejan los criterios antes mencionados. Como es
habitual, el entrenamiento y la evaluación de los modelos se hicieron sobre un
conjunto de entrenamiento y otro de \textit{testeo} respectivamente. Los
modelos evaluados fueron los siguientes: regresión logística, regresión
logísitca con penalización LASSO, bosques aleatorios, análisis discriminante
lineal y $k$ vecinos más cercanos. De estos, los dos que mostraron el mejor
desempeño fueron los bosques aleatorios y la regresión logística con
penalización LASSO con las variables originales normalizadas en ambos casos.

El informe se estructura de la siguiente manera. En la Sección \ref{sec:aed} se
da una descripción más detallada de los datos con los que se entrenaron los
modelos de clasificación y se resaltan algunos aspectos identificados que
tuvieron impacto sobre la estrategia de modelado. En la Sección
\ref{sec:estmod} se describe esta estrategia, en particular, algunas decisiones
generales de preprocesamiento de los datos, el método de validación, la
elección de las métricas de evaluación y su justificación y los modelos
específicos a implementar. En la Sección \ref{sec:res} se presentan los
resultados de la evaluación de todos los modelos implementados sobre los datos
de testeo y se detallan más específicamente los dos que mejor desempeño
tuvieron junto con el procedimiento completo de modelado. Finalmente, en la
Sección \ref{sec:con} se discuten brevemente los resultados y se resaltan los
puntos que resultaron más problemáticos en el desarrollo del trabajo práctico.
